% =====================================================================
%  CHAPTER 4: 儿童少年身体发育（对应大纲第四章）
%  内容来源：往届笔记第三章，参考PPT第三章
% =====================================================================
\chapter{儿童少年身体发育}
\setcounter{chapter}{4}
\chapterintro{
\textbf{掌握：}儿童少年体格发育、儿童少年体能发育、儿童少年体成分发育和脑发育。
}

% ===== 掌握内容 =====
\section{儿童少年体格发育}

\begin{defbox}
\textbf{体格发育（physical growth）：}人体外部形态、身体比例和体型等方面随年龄发生的变化。
\end{defbox}

\subsection{体格阶段性变化}

\begin{keybox}
\textbf{一、第一生长突增期：胎儿→2岁}

（1）\textbf{身高：}孕中期→1岁末；孕中期（4-6个月）增长量占胎儿期总增长量的50\%；出生后第1年增长约25cm。

（2）\textbf{体重：}孕晚期→1岁末；孕晚期（7-9个月）增长量占胎儿期总增长量的70\%。

\textbf{二、相对稳定期：2岁→青春期发育前}

身高：增长5-7cm/年，体重：增长2-3kg/年。

\textbf{三、第二生长突增期/青春期生长突增（adolescent growth spurt）：青春期}

（1）\textbf{身高：}突增高峰年增长10-14cm，男童>女童。

（2）\textbf{体重：}每年增长4-7kg，突增高峰年增长8-10kg。

\textbf{四、生长停滞期：青春后期开始}
\end{keybox}

\subsection{身体比例变化}

\begin{keybox}
\textbf{1. 反映横向身体比例指标：}反映身体各横截面指标的相互关系。

（1）\textbf{身高胸围指数}（胸围/身高）和\textbf{肩盆宽指数}（肩宽/盆宽）：只能衡量局部（身体上部或下部），实测值重复性差，后者个体差异只在青春期显现。

（2）\textbf{BMI：}既能反映身体胖瘦，又能确切反映不同群体、年代身体充实度变化。许多发达国家特别重视对BMI高百分位数变化的动态监测，将其作为制定肥胖防治策略和措施的重要依据。

\textbf{2. 反映纵向身体比例指标：}反映个体下肢长度和线性高度（身高）的比例关系。

（1）\textbf{身高坐高指数} = 坐高/身高。

（2）\textbf{下肢长指数：}下肢长指数I =（身高-坐高）/身高×100\%，下肢长指数II =（身高-坐高）/坐高×100\%。
\end{keybox}

\subsection{体型变化}

\begin{defbox}
\textbf{体型（somatotype）：}对人体形态的总体描述和评定，是身体各部位大小比例的形态特征总和，由遗传因素决定，也受人体对环境的适应和诸多环境因素的综合影响。

\textbf{分类：}

\begin{enumerate}[leftmargin=*]
    \item \textbf{内胚层体型（endomorphy）：}身体圆胖，消化器官发达。
    \item \textbf{中胚层体型（wesomorphy）：}身体健壮，骨骼肌肉发达。
    \item \textbf{外胚层体型（ectomorphy）：}身体瘦长，神经感觉系统发达。
\end{enumerate}
\end{defbox}

\section{儿童少年体能发育}

\begin{defbox}
\textbf{体能（physical fitness）/ 体适能：}人体具备的能胜任日常工作和学习而不感到疲劳，同时有余力能充分享受休闲娱乐生活，又可应付突发紧急情况的能力。
\end{defbox}

\subsection{体能分类}

\begin{keybox}
\textbf{1. 健康相关体能（health-related physical fitness）：}体能的基础部分，其目标是为身体健康和高质量，泛指人的体质，维持身体健康、提高工作、学习和生活效率所必需的基本技能。

指标体系：体成分、心肺耐力、柔韧性、肌力、肌耐力。是儿童少年为维持健康、提高生活质量等需求的追求。

\textbf{2. 运动技能相关体能（sports-related physical fitness）：}建立在健康相关体能基础上，属于较高体能需求层次，目标是为比赛取胜奖项。
\end{keybox}

\subsection{儿童少年体能发育特点}

\begin{keybox}
\textbf{1. 不均衡性：}不同体能指标在不同年龄发育速度有快有慢。

\textbf{2. 阶段性：}
\begin{enumerate}[leftmargin=*]
    \item \textbf{快速增长阶段：}男童6-14岁，女童6-12岁。
    \item \textbf{慢速增长阶段：}男童15-18岁，女童12-15岁。
    \item \textbf{恢复性生长：}仅女童16-18岁。
    \item \textbf{稳定阶段：}男性19-25岁，女性19-22岁。
\end{enumerate}

\textbf{3. 不平衡性}

\textbf{4. 性别特征}
\end{keybox}

\subsection{儿童少年体力活动}

\begin{defbox}
\textbf{体力活动（physical activity）：}WHO定义，由骨骼肌收缩产生需要能量消耗的任何身体运动，包括进行工作、家务、旅游、游戏、娱乐等。
\end{defbox}

\section{儿童少年体成分发育}

\begin{defbox}
\textbf{身体成分（body composition）/ 体成分：}人体总重量中不同身体成分的构成比例，属化学生长（chemical growth）范畴。
\end{defbox}

\subsection{体成分模型}

\begin{keybox}
\textbf{1. 二成分模型：}体脂重（FM）+ 去脂体重（FFM），主要使用。

\textbf{2. 多成分模型：}二成分模型的细化。
\end{keybox}

\subsection{体成分发育的性别特征}

\begin{keybox}
\textbf{1. 体成分发育的性别差异}

\begin{enumerate}[leftmargin=*]
    \item 生命早期，男女体脂量接近。
    \item 儿童期，男、女体脂均下降。
    \item 青春期，女童变化以体脂增加为主，男童体脂有下降。
    \item 体脂分布的性别差异在青春期前即有表现；男童上半身的皮下和内脏脂肪蓄积，女童皮下脂肪蓄积在臀部。
\end{enumerate}

\textbf{2. 瘦体重发育的性别特征}

\begin{enumerate}[leftmargin=*]
    \item 生命早期开始，男孩就显出更高的瘦体重，故其体重>同年龄女孩。
    \item 儿童期，男孩仍较女孩有较多瘦体重增长。
    \item 进入青春期后，男孩主要表现为瘦体重增加，女孩的瘦体重增量仅为男孩的一半。
    \item 随着近年来性早熟（女性更多）现象增加，瘦体重的性别差异出现时间明显提前。
\end{enumerate}
\end{keybox}

\section{脑发育}

\begin{keybox}
\textbf{生命早期脑发育：}

生命早期（受精卵形成-出生后2年），大脑以惊人的速度生长。胎儿期的最后3个月（孕28周）和婴儿出生后2年被称作"大脑发育加速期"，因为成人脑一半以上（75\%）的重量都在该阶段获得。
\end{keybox}

% ----- 复习题 -----
\reviewcontent{
\section{复习题}

\begin{enumerate}[leftmargin=*, itemsep=8pt]
    \item \textbf{【名词解释】}体格发育；体型；体能/体适能；健康相关体能；体力活动；身体成分/体成分
    \item \textbf{【简答题】}简述儿童少年体格发育的四个阶段及其特点。
    \item \textbf{【简答题】}简述反映横向和纵向身体比例的指标及其意义。
    \item \textbf{【简答题】}简述三种体型的分类及特点。
    \item \textbf{【简答题】}简述儿童少年体能发育的阶段性特点。
    \item \textbf{【简答题】}简述体成分发育的性别特征。
    \item \textbf{【简答题】}简述生命早期脑发育的特点。
\end{enumerate}

\subsection*{参考答案要点}

\begin{enumerate}[leftmargin=*, itemsep=5pt]
    \item \textbf{体格发育}：人体外部形态、身体比例和体型等方面随年龄发生的变化。\textbf{体型}：对人体形态的总体描述和评定，是身体各部位大小比例的形态特征总和。\textbf{体能/体适能}：人体具备的能胜任日常工作和学习而不感到疲劳，同时有余力充分享受休闲娱乐生活，又可应付突发紧急情况的能力。\textbf{健康相关体能}：体能的基础部分，维持身体健康、提高工作、学习和生活效率所必需的基本技能。\textbf{体力活动}：由骨骼肌收缩产生需要能量消耗的任何身体运动。\textbf{身体成分/体成分}：人体总重量中不同身体成分的构成比例，属化学生长范畴。

    \item (1)第一生长突增期（胎儿→2岁）：身高孕中期占胎儿期50\%，出生后第1年增长约25cm，体重孕晚期占70\%；(2)相对稳定期（2岁→青春期前）：身高5-7cm/年，体重2-3kg/年；(3)第二生长突增期（青春期）：身高10-14cm/年，体重4-7kg/年（高峰年8-10kg）；(4)生长停滞期（青春后期开始）。

    \item 横向：身高胸围指数和肩盆宽指数（衡量局部）；BMI（反映胖瘦和身体充实度变化，是制定肥胖防治策略的重要依据）。纵向：身高坐高指数=坐高/身高；下肢长指数I=（身高-坐高）/身高×100\%。

    \item 内胚层体型：身体圆胖，消化器官发达；中胚层体型：身体健壮，骨骼肌肉发达；外胚层体型：身体瘦长，神经感觉系统发达。

    \item (1)快速增长阶段：男6-14岁，女6-12岁；(2)慢速增长阶段：男15-18岁，女12-15岁；(3)恢复性生长：仅女16-18岁；(4)稳定阶段：男19-25岁，女19-22岁。同时具有不均衡性、不平衡性和性别特征。

    \item 体成分：(1)生命早期男女接近→儿童期均下降→青春期女以体脂增加为主、男体脂下降；(2)体脂分布差异在青春期前即有表现，男上身上皮下和内脏脂肪蓄积，女臀部皮下脂肪蓄积。瘦体重：(1)生命早期男孩即更高；(2)儿童期男孩增长更多；(3)青春期男孩主要表现为瘦体重增加，女孩增量仅为男孩一半；(4)性早熟增加使性别差异出现时间提前。

    \item 生命早期（受精卵形成-出生后2年）大脑以惊人速度生长。胎儿期最后3个月（孕28周）和婴儿出生后2年为"大脑发育加速期"，成人脑75\%的重量在该阶段获得。
\end{enumerate}
}

\newpage
